\documentclass[11pt,a4paper]{article}
\usepackage[utf8]{inputenc}
\usepackage[T1]{fontenc}
\usepackage[margin=2.2cm]{geometry}
\usepackage{amsmath,amssymb}
\usepackage{booktabs}
\usepackage{xcolor}
\usepackage{fontawesome5}
\usepackage{hyperref}
\usepackage{tcolorbox}
\tcbuselibrary{skins}

\definecolor{BioNavy}{RGB}{26, 54, 93}
\definecolor{BioTeal}{RGB}{13, 148, 136}
\definecolor{BioAmber}{RGB}{217, 119, 6}

\hypersetup{
    colorlinks=true,
    linkcolor=BioNavy,
    urlcolor=BioTeal
}

\title{\textbf{\LARGE Guide Pratique d'Installation \& Prise en Main}\\
\Large Python, Anaconda \& Jupyter Notebook pour les Biologistes}
\author{\textbf{Dr. Sarra BENMOUMOU (Ph.D.)}\\
\small Faculté des Sciences --- Département de Biologie\\
\small Master 2 Biochimie Appliquée --- Université M'Hamed Bougara de Boumerdès (UMBB)}
\date{Année Universitaire 2024 -- 2025}

\begin{document}

\maketitle

\vspace{-5mm}
\begin{center}
  \rule{\textwidth}{1pt}
\end{center}

\section*{\faGraduationCap\hspace{0.5em}Bienvenue dans le Calcul Scientifique Biologique !}
Chers étudiants,\\
En tant que futurs chercheurs ou cadres en biochimie et biotechnologie, vous manipulerez des données de plus en plus complexes (cinétiques enzymatiques, spectres, ELISA, profilages métabolomiques).

\textbf{Rassurez-vous :} vous n'avez besoin d'\textbf{aucune formation préalable en informatique}. Python et les carnets interactifs \textbf{Jupyter Notebook} ont été choisis pour leur extrême simplicité, leur lisibilité et leur capacité à générer en quelques lignes des graphiques prêts pour la publication scientifique.

\section*{\faDownload\hspace{0.5em}Étape 1 : Téléchargement Gratuit d'Anaconda}
\textbf{Anaconda} est la distribution scientifique de référence dans le monde. Elle installe en un seul paquet :
\begin{itemize}
  \item Python 3.x et l'interface conviviale \textbf{Anaconda Navigator}.
  \item L'environnement interactif de travail \textbf{Jupyter Notebook}.
  \item Toutes les bibliothèques utiles en biochimie : \texttt{pandas}, \texttt{scipy}, \texttt{matplotlib}, \texttt{seaborn}, \texttt{scikit-learn}.
\end{itemize}

\textbf{Lien officiel de téléchargement :} \url{https://www.anaconda.com/download}

\section*{\faCogs\hspace{0.5em}Étape 2 : Installation Pas-à-Pas}
\begin{itemize}
  \item \textbf{Sous Windows :} Double-cliquez sur le fichier \texttt{.exe} téléchargé. Cliquez sur \textit{Next}, acceptez la licence (\textit{I Agree}), choisissez \textit{Just Me (recommended)} et conservez tous les paramètres par défaut.
  \item \textbf{Sous macOS :} Ouvrez le fichier \texttt{.pkg} et suivez les instructions à l'écran.
  \item \textbf{Sous Linux :} Ouvrez un terminal et tapez \texttt{bash Anaconda3-xxxx-Linux-x86\_64.sh}.
\end{itemize}

\section*{\faRocket\hspace{0.5em}Étape 3 : Lancement de Jupyter Notebook}
\begin{enumerate}
  \item Ouvrez le menu Démarrer de votre ordinateur et lancez \textbf{Anaconda Navigator}.
  \item Dans la mosaïque d'applications, repérez \textbf{Jupyter Notebook} et cliquez sur le bouton vert \textbf{Launch}.
  \item Votre navigateur internet s'ouvre : vous êtes prêt à travailler !
\end{enumerate}

\section*{\faKeyboard\hspace{0.5em}Étape 4 : Les Raccourcis Indispensables}
\begin{center}
\begin{tabular}{ll}
  \toprule
  \textbf{Action Souhaitée} & \textbf{Raccourci Clavier} \\
  \midrule
  \textbf{Exécuter le calcul de la cellule sélectionnée} & \texttt{Maj (Shift) + Entrée} \\
  Créer une nouvelle cellule en dessous & \texttt{Échap} puis \texttt{B} \\
  Transformer une cellule en texte explicatif (Markdown) & \texttt{Échap} puis \texttt{M} \\
  Transformer une cellule en code de calcul & \texttt{Échap} puis \texttt{Y} \\
  Supprimer une cellule inutile & \texttt{Échap} puis \texttt{D, D} \\
  \bottomrule
\end{tabular}
\end{center}

\vspace{5mm}
\begin{tcolorbox}[colback=BioTeal!10, colframe=BioTeal, title={\faFlask\hspace{0.5em}Exemple Immédiat : Tracer une Gamme Étalon en 3 Lignes}]
\small
\begin{verbatim}
import pandas as pd
import seaborn as sns
import matplotlib.pyplot as plt

# Chargement des mesures
df = pd.read_csv('../datasets/dosage_bradford.csv')

# Tracé automatique avec droite de régression
sns.lmplot(data=df, x='BSA_ug_mL', y='Absorbance_595nm', color='#1A365D')
plt.title("Gamme Étalon Bradford (BSA à 595 nm)")
plt.show()
\end{verbatim}
\end{tcolorbox}

\vspace{5mm}
\begin{center}
  \footnotesize
  \color{gray}
  Module de Biostatistiques --- Responsable pédagogique : Dr. Sarra BENMOUMOU (Ph.D.) --- UMBB
\end{center}

\end{document}
